
\subsubsection{Model-Based Code Evaluation}

Traditional code evaluation relies on execution-based metrics such as pass@k \citep{chen2021evaluating}, which measure functional correctness through test execution. Model-based alternatives emerged to approximate execution signals. CodeBERTScore \citep{zhou2023codebertscore} computes semantic similarity using code-pretrained embeddings. Naik (2024) reports that CodeBERTScore has 0.16 correlation with functional correctness, indicating limited reliability for correctness judgment.

\subsubsection{LLM-as-Judge for Code}

Moon et al. (2025) identify six bias types in LLM code judges: sensitivity to variable names, comments, formatting, and other superficial features. Crupi et al. (2025) evaluate eight LLMs as judges on Java and Python methods, finding that GPT-4-turbo performs best but ''frequently misjudges'' correctness. Critically, their study does not control for scale. Wang et al. (2025) introduce MCTS-Judge, applying test-time compute to improve single-model accuracy from 41\% to 80\%.

\subsubsection{Ensemble Methods}

Ensemble methods assume component errors are approximately IID. Shu (2026) demonstrates that LLM judge panels share fundamental error modes even when architecturally diverse. The present findings extend this observation: within scale-diverse panels, errors are not only correlated but systematically asymmetric, making majority voting counterproductive.

